\section{Experimental Design}

\subsection{Factorial confirmatory comparison}
\begin{draftpoints}
  \item main MNIST design crosses architecture update type (fixed CL versus growing NDL) with replay regime (none, original data, or intrinsic replay)
  \item Original-data replay is an oracle upper bound because it reopens historical images; intrinsic replay is the paper-shaped memory condition because it stores only latent class statistics
  \item Ten independent seeds, 42--51, are used for the full MNIST confirmation; two-sided Student-$t$ confidence intervals summarize seed variability
  \item principal comparisons are paired within replay regime: NDL versus its exact capacity-matched CL control
\end{draftpoints}

\begin{table}[htbp]
\centering
\caption{Complete experiment matrix and evidential role of each family.}
\label{tab:experiment-matrix}
\small
\begin{tabularx}{\textwidth}{p{2.7cm}p{1.4cm}p{2.1cm}p{2.4cm}p{1.5cm}Y}
\toprule
Condition & Growth & Replay & Capacity relation & Seeds & Purpose \\
\midrule
CL, no replay & No & None & Matched to NDL seed & 42--51 & Fixed learner without memory \\
NDL, no replay & Yes & None & Organic endpoint & 42--51 & Isolate growth without memory \\
CL + original data & No & Old images & Matched to NDL & 42--51 & Clean replay upper bound \\
NDL + original data & Yes & Old images & Organic endpoint & 42--51 & Growth with ideal replay \\
CL + intrinsic replay & No & Latent Gaussian & Matched to NDL & 42--51 & Isolate generated replay \\
NDL + intrinsic replay & Yes & Latent Gaussian & Organic endpoint & 42--51 & Central paper-shaped condition \\
SD-19 CL / NDL & No / Yes & Original data & Paired endpoints & 42--44 & Larger-dataset feasibility \\
Organic-growth screens & Yes & Original data & Multiple policies & seed 42 & Test whether demand, not caps, determines width \\
Predictive coding & No & None / original & Fixed endpoint & 42--44 & Post-replication optimizer extension \\
Optimized NDL / CL & Yes / No & Original data & Exact endpoint & 45--47 & Post-replication global-coupling extension \\
\bottomrule
\end{tabularx}
\end{table}

\subsection{Outcomes}
\begin{draftpoints}
  \item Macro MSE gives equal weight to each learned class and is the primary endpoint because class frequencies should not dominate the continual-learning comparison
  \item Foreground-pixel MSE emphasizes reconstruction of digit or character strokes rather than the mostly black background
  \item Mean positive forgetting averages only increases in old-class error, separating retention failures from improvements that can make signed averages misleading
  \item Structural outcomes include final encoder widths, nodes added per layer, quota-stop fraction, updated-class fraction, and unresolved exhausted levels
  \item Efficiency outcomes include parameter count, pretraining and incremental update counts, wall time, peak CUDA allocation and reservation, and peak process resident memory
\end{draftpoints}

\subsection{Class-level analysis}
\begin{draftpoints}
  \item Per-class results separate the initially learned digits 1 and 7 from early and late incoming classes
  \item Incoming-class error measures plasticity, while changes in earlier classes measure stability; both are required because a method can retain old classes by failing to learn the new class
  \item preserved per-class dataset contains 600 MNIST seed observations and 96 SD-19 seed observations, enabling class-specific comparison without relying on a single aggregate score
\end{draftpoints}

\subsection{Confirmatory versus exploratory evidence}
\begin{draftpoints}
  \item Paper-compatible parameter screens are used to resolve undocumented implementation choices before the frozen full-curriculum confirmation
  \item Global reconstruction coupling, predictive coding, and usage-dependent plasticity change the paper mechanism and therefore cannot rescue the original replication claim even if they improve performance
  \item SD-19 is described as a feasibility comparison because only three seeds and a restricted follow-up are preserved, not the complete twenty-order curriculum of the original paper
\end{draftpoints}

\begin{figure}[htbp]
\centering
\resizebox{\textwidth}{!}{%
\begin{tikzpicture}[
node distance=7mm and 8mm,
stage/.style={draw,rounded corners,minimum width=30mm,minimum height=10mm,align=center,font=\scriptsize},
gate/.style={draw,diamond,aspect=2.2,align=center,font=\scriptsize,inner sep=1pt},
arrow/.style={-{Latex[length=2mm]},thick}
]
\node[stage,fill=ProtocolBlue] (audit) {Fidelity audit\\and defect fixes};
\node[stage,right=of audit] (base) {Base-model\\competence};
\node[stage,right=of base] (screen) {Growth and\\protocol screens};
\node[gate,right=of screen] (gate) {Acquisition, retention,\\and organicity gates};
\node[stage,right=of gate,fill=ProtocolLightPurple] (confirm) {Ten-seed frozen\\confirmation};
\node[stage,below=of confirm,fill=ProtocolRed] (matched) {Exact endpoint-matched\\controls};
\node[stage,left=of matched] (extension) {Non-paper\\extensions};
\draw[arrow] (audit) -- (base);
\draw[arrow] (base) -- (screen);
\draw[arrow] (screen) -- (gate);
\draw[arrow] (gate) -- node[above,font=\scriptsize]{pass} (confirm);
\draw[arrow] (confirm) -- (matched);
\draw[arrow] (matched) -- (extension);
\draw[arrow,dashed] (gate.south) |- node[pos=.25,right,font=\scriptsize]{mechanistic follow-up} (extension.east);
\end{tikzpicture}
}
\caption{Selection and evidence flow. Screening informs frozen confirmatory settings; post-replication extensions remain outside the confirmatory claim.}
\label{fig:selection-flow}
\end{figure}
